This case demonstrates that, when tuned with Flow-GRPO, the \textit{Action Planner} first attempts several solutions, recognizes their ineffectiveness, resorts to a brute-force approach, and finally verifies the result using a search engine.

\begin{custombox}[title=Example from \# GameOf24]{toolcardborder}
\textbf{Question:} Using the numbers [1, 1, 6, 9], create an expression that equals 24. You must use basic arithmetic operations (+, -, *, /) and parentheses.

\vspace{1em}

\textbf{Answer:} (1+1)*9+6
\end{custombox}

\begin{small}
\begin{center}

\begin{casebox}{successbg}{successframe}{Success Case: \model w/ Flow-GRPO}
    \textbf{Action Step 1: Initial Attempt}
    \begin{itemize}[nosep, leftmargin=*]
        \item \textbf{Tool:} \smalltt{Python Coder}
        \item \textbf{Sub-goal:} Generate and test all possible arithmetic expressions using the numbers [1, 1, 6, 9] exactly once to check if any combination equals 24.
        \item \textbf{Command:} 
        
        \smalltt{tool.execute(query="Using the numbers [1, 1, 6, 9], create an expression that equals 24 using basic arithmetic operations and parentheses.")}
        \item \textbf{Result:}
        \begin{casecode}{successbg}
\# Given numbers

numbers = [1, 1, 6, 9]

\# Calculate the expression

result = (numbers[0] + numbers[1]) * (numbers[3] - numbers[2])

\# Print the result

print(f"The result of the expression ({numbers[0]} + {numbers[1]} ) * ({numbers[3]} - {numbers[2]}) is: {result}")
        \end{casecode}
        The result of the expression (1 + 1 ) * (9 - 6) is: 6
    \end{itemize}

    \vspace{1em}

    \textbf{Action Step 2: Tool Error}
    \begin{itemize}[nosep, leftmargin=*]
        \item \textbf{Tool:} \smalltt{Python Coder}
        \item \textbf{Sub-goal:} Generate and test all possible permutations of the numbers with different operator placements to find an expression that equals 24.
        \item \textbf{Command:} 
        
        \smalltt{tool.execute(query="Using the numbers [1, 1, 6, 9], create an expression with +,-,*,/ that equals 24")}
        \item \textbf{Result:} \textcolor{orange}{Error: eval() arg 1 must be a string, bytes or code object}
    \end{itemize}

    \vspace{1em}

    \textbf{Action Step 3: Second Attempt}
    \begin{itemize}[nosep, leftmargin=*]
        \item \textbf{Tool:} \smalltt{Python Coder}
        \item \textbf{Sub-goal:} Generate all possible combinations of the four numbers with every permutation of operators (+,-,*,/) to find a valid equation resulting in exactly 24.
        \item \textbf{Command:} 
        
        \smalltt{tool.execute(query="Using the numbers [1, 1, 6, 9] create an expression with +,-,*,/ that equals to '24'")}
        \item \textbf{Result:}
        \begin{casecode}{successbg}
\# Given numbers

numbers = [1, 1, 6, 9]
        \end{casecode}
        \begin{casecode}{successbg}
\# Calculate the expression

result = (numbers[2] - numbers[0]) * (numbers[3] / numbers[1]) / numbers[1]

\# Print the result

print(f"The result of the expression is: {result}")
        \end{casecode}
        The result of the expression is: 45.0
    \end{itemize}

    \vspace{1em}

    \textbf{Action Step 4: Systematic Search (Breakthrough)}
    \begin{itemize}[nosep, leftmargin=*]
        \item \textbf{Tool:} \smalltt{Python Coder}
        \item \textbf{Sub-goal:} Generate and evaluate all possible permutations of the four numbers with every permutation of operators (+,-,*,/) to find a valid equation resulting in exactly 24.
        \item \textbf{Command:} 
        
        \smalltt{tool.execute(query="Using the numbers [1, 1, 6, 9] create an expression that equals 24 using +,-,/,*")}
        \item \textbf{Result:} \textbf{Successfully found valid expressions:} The expression (1+1)*9+6 equals 24. The expression 6+(1+1)*9 equals 24. The expression 6+(9*(1+1)) equals 24. The expression 9*(1+1)+6 equals 24.
    \end{itemize}

    \vspace{1em}

    \textbf{Action Step 5: Confused Verification}
    \begin{itemize}[nosep, leftmargin=*]
        \item \textbf{Tool:} \smalltt{Python Coder}
        \item \textbf{Sub-goal:} Simplify and verify if any minor adjustments to the existing promising expression can achieve an exact match of 24.
        \item \textbf{Command:} 
        
        \smalltt{tool.execute(query="Using the numbers [1, 1, 6, 9] and basic arithmetic operations create an expression that equals 24: (9 * (6 * (1 + 1)))")}
        \item \textbf{Result:} \textcolor{orange}{The result of the expression (9 * (1 + 1) * 6) is: 108}
    \end{itemize}

    \vspace{1em}

    \textbf{Action Step 6: Final Verification}
    \begin{itemize}[nosep, leftmargin=*]
        \item \textbf{Tool:} \smalltt{Google Search}
        \item \textbf{Sub-goal:} Search online for examples or discussions related to creating expressions using given numbers that equal a target value.
        \item \textbf{Command:} 
        
        \smalltt{tool.execute(query="using the numbers [1, 1, 6, 9] create an expression that equals 24", add\_citations=True)}
        \item \textbf{Result:} Here's an expression using the numbers that equals 24: (1 + 1) * 9 + 6 = 24
    \end{itemize}

    \vspace{1em}

    \textbf{Execution Verifier:} \textcolor{successframe}{PASS, STOP (via exhaustive search and verification)}

    \vspace{1em}

    \textbf{Solution Generator: } \textcolor{successframe}{After a multi-step process of trial, error, and verification, the final correct expression is \textbf{(1 + 1) * 9 + 6 = 24}.}

\end{casebox}

\begin{casebox}{failbg}{failframe}{Failure Case: \model w/o Flow-GRPO}
    \textbf{Action Step 1:}
    \begin{itemize}[nosep, leftmargin=*]
        \item \textbf{Tool:} \smalltt{Python Coder}
        \item \textbf{Sub-goal:} Generate and test Python code snippets to find an arithmetic expression using the numbers [1, 1, 6, 9] exactly once that equals 24.
        \item \textbf{Command:} 
        
        \smalltt{tool.execute(query="Using the numbers [1, 1, 6, 9], create an expression that equals 24 using basic arithmetic operations and parentheses.")}
        \item \textbf{Result:}
        \begin{casecode}{failbg}
\# Define the numbers

numbers = [1, 1, 6, 9]

\# Calculate the expression

result = (6 * 9) - ((1 + 1) * 15)

\# Print the result with a descriptive message

print("The result of the expression is:", result)
        \end{casecode}
        The result of the expression is: 24
    \end{itemize}

    \vspace{1em}

    \textbf{Execution Verifier:} \textcolor{red}{STOP}

    \vspace{1em}

    \textbf{Solution Generator: } \textcolor{red}{The agent incorrectly concludes the answer is $(6 \times 9) - ((1 + 1) \times 15)$, failing to notice it used the number \textbf{15}, which was not in the allowed input set.}
\end{casebox}


\end{center}
\end{small}
